\documentclass[a4paper,12pt]{article}
\usepackage[T2A]{fontenc}
\usepackage[utf8]{inputenc}
\usepackage[russian]{babel}
\usepackage{amsmath,amssymb,mathtools}
\usepackage{helvet}
\renewcommand{\familydefault}{\sfdefault}
\usepackage{icomma}
\usepackage[left=22mm,right=22mm,top=23mm,bottom=25mm]{geometry}
\usepackage{microtype}
\usepackage{booktabs,tabularx,array}
\usepackage{enumitem}
\usepackage{xcolor}
\usepackage{tikz}
\usepackage{hyperref}
\usepackage{fancyhdr}
\usepackage{parskip}
\definecolor{accent}{HTML}{006F72}
\definecolor{darkaccent}{HTML}{004E53}
\definecolor{ink}{HTML}{202C33}
\definecolor{muted}{HTML}{5A676B}
\definecolor{errorred}{HTML}{AE2933}
\hypersetup{colorlinks=true,linkcolor=darkaccent,pdfauthor={Лёвин Артём Александрович},pdftitle={Ревью домашней работы — 19 ноября 2026 г.}}
\setlength{\parindent}{0pt}
\setlength{\parskip}{0.6em}
\setlist[enumerate]{leftmargin=*,itemsep=0.35em,topsep=0.45em}
\pagestyle{fancy}
\fancyhf{}
\renewcommand{\headrulewidth}{0pt}
\renewcommand{\footrulewidth}{0.3pt}
\fancyfoot[C]{\footnotesize\color{muted}Лёвин Артём Александрович эксклюзивно для Ксении Клыковой}
\fancyfoot[R]{\footnotesize\color{muted}\thepage}
\newcommand{\subblock}[1]{\par\medskip{\large\bfseries\color{darkaccent}#1}\par\smallskip}
\newcommand{\errlabel}{\textcolor{errorred}{\bfseries Ошибка:}\,}
\newcommand{\idea}{\textcolor{darkaccent}{\bfseries Ключевая идея:}\,}
\newcommand{\result}{\textcolor{accent}{\bfseries Итог:}\,}
\newcommand{\problem}[2]{\phantomsection\label{task:#1}\section*{Задание №\,#1. #2}}
\begin{document}
\thispagestyle{empty}
\begin{center}
\vspace*{45mm}
{\large\color{accent}\textbf{МАТЕМАТИКА}}\par
\vspace{17mm}
{\fontsize{26}{31}\selectfont\bfseries\color{ink}Ревью домашней работы}\par
\vspace{8mm}
{\large\color{muted}Профильная математика \textbullet\ Разбор ошибок}\par
\vspace{22mm}
{\normalsize 19 ноября 2026 г.}\par
\vspace{6mm}
{\normalsize Автор: Лёвин Артём Александрович}\par
\vfill
\begin{tikzpicture}[x=1cm,y=1cm]
\draw[accent!70,line width=0.9pt] (-5.7,0) .. controls (-3.8,0.65) and (-1.6,-0.58) .. (0,0)
  .. controls (1.6,0.58) and (3.8,-0.65) .. (5.7,0);
\end{tikzpicture}
\vspace{11mm}
\end{center}
\clearpage

\section*{Сводная таблица}
В таблицу включены только задания, в которых обнаружено нарушение логики решения или решение не доведено до ответа. Номера ведут к подробным разборам.
\vspace{4mm}
\renewcommand{\arraystretch}{1.48}
\setlength{\tabcolsep}{4pt}
\small
\begin{tabularx}{\linewidth}{@{}p{11mm}p{23mm}X p{24mm}p{29mm}@{}}
\toprule
\textbf{№} & \textbf{Тема} & \textbf{Суть ошибки} & \textbf{Ваш ответ} & \textbf{Правильный ответ} \\
\midrule
\hyperref[task:12]{12} & Дробно-рациональная функция & Потеря минуса в одной из промежуточных дробей. & $-3{,}2$ & $-3{,}2$ \\
\hyperref[task:13]{13} & Дробно-рациональное уравнение & Ошибка в знаке при раскрытии скобок; затем потерян отдельный член уравнения. Решение не завершено. & Не указан & $x=-5,\;\dfrac{-1-\sqrt{65}}{2}$ на заданном отрезке \\
\bottomrule
\end{tabularx}
\normalsize
\vspace{10mm}
\textbf{На что обратить внимание.} В первом случае достаточно восстановить корректную цепочку равенств и проверить знак результата. Во втором случае полезнее упростить структуру уравнения с помощью одной замены, чем раскрывать все скобки сразу.
\clearpage

\thispagestyle{empty}
\vspace*{0.36\textheight}
\begin{center}
{\Huge\bfseries\color{darkaccent}Разбор ошибок}\par\vspace{5mm}
{\normalsize\color{muted}Точные места ошибок, исправления и проверка}\par
\end{center}
\clearpage

\problem{12}{Знак при вычислении значения функции}
\subblock{Текст задания}
На рисунке задана функция
\[
 f(x)=\frac{k}{x+a}.
\]
Её график проходит через точки $A(3;-1)$ и $B(1;-2)$. Найдите $f(0{,}25)$.

\subblock{Ваше решение}
Из координат точек правильно составлена система
\[
-1=\frac{k}{3+a},\qquad -2=\frac{k}{1+a}.
\]
Получены $a=1$ и $k=-4$. При подстановке найденных параметров появилась запись, в которой положительная дробь $\dfrac{4}{1{,}25}$ приравнена к отрицательному результату $-3{,}2$.

\subblock{Ваш ответ}
$-3{,}2$.

\subblock{Ошибка}
\errlabel потеря знака минус в промежуточном преобразовании. Конечный ответ записан верно, но одно из равенств в вычислении неверно.

\subblock{Где именно возникает ошибка}
Последний правильный шаг:
\[
f(0{,}25)=\frac{-4}{0{,}25+1}=\frac{-4}{1{,}25}.
\]
В следующем переходе минус у числителя исчезает:
\[
\frac{-4}{1{,}25}\ \textcolor{errorred}{\ne}\ \frac{4}{1{,}25}.
\]

\subblock{Почему это неверно}
Числитель $-4$ отрицателен, знаменатель $1{,}25$ положителен. По правилу знаков при делении частное отрицательного числа на положительное отрицательно. Исчезновение минуса изменяет значение дроби на противоположное. Два соседних звена цепочки равенств должны обозначать одно и то же число. Правильный ответ в последней строке не исправляет неверное промежуточное равенство.

\subblock{Ключевая идея}
\idea перед вычислением десятичной дроби определите знак выражения по знакам числителя и знаменателя. Для функции $k/(x+a)$ отдельно контролируйте, что $x+a\ne0$.

\subblock{Исправление}
Продолжаем от последнего правильного шага, сохраняя минус:
\[
\frac{-4}{1{,}25}=-\frac{4}{1{,}25}=-3{,}2.
\]

\subblock{Полное правильное решение}
Подставим координаты обеих точек в формулу функции:
\[
\begin{cases}
-1=\dfrac{k}{3+a},\\[5pt]
-2=\dfrac{k}{1+a}.
\end{cases}
\]
Умножим первое равенство на $3+a$, второе --- на $1+a$:
\[
\begin{cases}
k=-(3+a),\\
k=-2(1+a).
\end{cases}
\]
Приравняем два выражения для $k$:
\[
-3-a=-2-2a
\quad\Longrightarrow\quad a=1.
\]
Тогда $k=-2(1+1)=-4$ и
\[
f(x)=\frac{-4}{x+1},\qquad x\ne-1.
\]
Подставим $x=0{,}25$:
\[
f(0{,}25)=\frac{-4}{0{,}25+1}=
\frac{-4}{1{,}25}=\frac{-400}{125}=-\frac{16}{5}=-3{,}2.
\]

\subblock{Проверка}
Проверим найденную функцию в двух известных точках:
\[
f(3)=\frac{-4}{4}=-1,\qquad
f(1)=\frac{-4}{2}=-2.
\]
Обе координаты подтверждены. При $x=0{,}25$ знаменатель положителен, числитель отрицателен; знак $f(0{,}25)$ действительно отрицательный.

\subblock{Итог}
\result $\boxed{f(0{,}25)=-3{,}2}$. Главный навык для закрепления --- следить за знаком дроби на каждом шаге.

\subblock{Задача на закрепление}
График функции $g(x)=\dfrac{m}{x+b}$ проходит через точки $A(2;-1)$ и $B(0;-3)$. Найдите $g(0{,}5)$.
\clearpage

\problem{13}{Преобразование дробно-рационального уравнения}
\subblock{Текст задания}
\textbf{а)} Решите уравнение
\[
\frac{(x+3)^2}{5}+\frac{20}{(x+3)^2}
=8\left(\frac{x+3}{5}-\frac{2}{x+3}\right)+1.
\]
\textbf{б)} Найдите его корни, принадлежащие отрезку $[-6;-4]$.

\subblock{Ваше решение}
Вначале Вы отдельно приводите к общему знаменателю выражение в скобках справа и сумму дробей слева. Идея привести дроби к общему знаменателю сама по себе допустима. Однако в дальнейшей цепочке возникает неверный знак при раскрытии скобок: выражение $x^2+6x+9-10$ переписывается с отрицательным коэффициентом при $x$. После умножения скобок на $8$ в последующих строках также теряется отдельное слагаемое $+1$ из правой части исходного уравнения. В конце предпринята попытка исследовать числитель и знаменатель как самостоятельные уравнения; она не устанавливает равенство исходных частей. Надёжного окончательного решения и отбора корней на отрезке в записи нет.

\subblock{Ваш ответ}
Не указан.

\subblock{Ошибка}
\errlabel нарушение равносильности при преобразованиях уравнения: ошибка в знаке, пропуск слагаемого и переход к отдельным выражениям вместо сохранения исходного равенства.

\subblock{Где именно возникает ошибка}
Последний надёжно правильный шаг при приведении выражения в скобках к общему знаменателю:
\[
\frac{x+3}{5}-\frac{2}{x+3}
=\frac{(x+3)^2-10}{5(x+3)}
=\frac{x^2+6x-1}{5(x+3)}.
\]
Первый неверный шаг в рукописной цепочке --- замена $+6x$ на $-6x$ при упрощении числителя. Поэтому появившееся далее выражение с $8x^2-48x-8$ в числителе тоже неверно: при умножении на $8$ должно быть $8x^2+48x-8$. Кроме того, слагаемое $+1$ после скобок обязательно сохраняется.

\subblock{Почему это неверно}
При возведении суммы в квадрат действует формула
\[
(a+b)^2=a^2+2ab+b^2.
\]
Поэтому $(x+3)^2=x^2+6x+9$: средний член положителен. Последующее вычитание $10$ изменяет только свободный член, поскольку $9-10=-1$. Кроме того, каждое действие с уравнением должно сохранять обе части целиком: при замене скобки выражением или при приведении дробей к общему знаменателю внешний член $+1$ остаётся на месте. Наконец, равенство двух дробно-рациональных выражений нельзя заменить поиском нулей их отдельных числителей и знаменателей. Знаменатель при допустимых значениях переменной отличен от нуля.

\subblock{Ключевая идея}
\idea обратите внимание на повторяющуюся разность
\[
t=\frac{x+3}{5}-\frac{2}{x+3}.
\]
Левая часть уравнения выражается через квадрат этой разности. Такая замена избавляет от громоздкого раскрытия четвёртых степеней.

\subblock{Исправление}
Сначала выпишем область допустимых значений: $x+3\ne0$, следовательно $x\ne-3$. Пусть
\[
u=\frac{x+3}{5},\qquad v=\frac{2}{x+3}.
\]
Тогда $uv=\dfrac25$, а исходное уравнение принимает вид
\[
5u^2+5v^2=8(u-v)+1.
\]
По формуле $u^2+v^2=(u-v)^2+2uv$ получаем
\[
5(u-v)^2+5\cdot\frac45=8(u-v)+1,
\]
то есть
\[
5(u-v)^2+4=8(u-v)+1.
\]
Остаётся обозначить $t=u-v$ и решить квадратное уравнение.

\subblock{Полное правильное решение}
\textbf{Шаг 1. Область допустимых значений.} Знаменатели содержат $x+3$ и $(x+3)^2$, поэтому
\[
\boxed{x\ne-3}.
\]

\textbf{Шаг 2. Замена.} Обозначим
\[
t=\frac{x+3}{5}-\frac{2}{x+3}.
\]
Вычислим квадрат:
\[
\begin{aligned}
t^2&=\frac{(x+3)^2}{25}
-2\cdot\frac{x+3}{5}\cdot\frac{2}{x+3}
+\frac{4}{(x+3)^2}\\[2pt]
&=\frac{(x+3)^2}{25}-\frac45+\frac{4}{(x+3)^2}.
\end{aligned}
\]
Умножая на $5$ и прибавляя $4$, получим точную левую часть исходного уравнения:
\[
5t^2+4=\frac{(x+3)^2}{5}+\frac{20}{(x+3)^2}.
\]
Следовательно, исходное уравнение равносильно
\[
5t^2+4=8t+1.
\]

\textbf{Шаг 3. Решаем уравнение с новой переменной.}
\[
5t^2-8t+3=0.
\]
Дискриминант:
\[
D=(-8)^2-4\cdot5\cdot3=64-60=4.
\]
По формуле корней
\[
t=\frac{8\pm2}{10},\qquad
\boxed{t=1\ \text{или}\ t=\frac35}.
\]

\textbf{Шаг 4. Возвращаемся к переменной $x$.}

\textit{Случай I: $t=1$.}
\[
\frac{x+3}{5}-\frac{2}{x+3}=1.
\]
Умножим на $5(x+3)$, что допустимо при $x\ne-3$:
\[
(x+3)^2-10=5(x+3).
\]
Раскрываем скобки, переносим все слагаемые влево:
\[
x^2+6x+9-10-5x-15=0,
\]
\[
x^2+x-16=0.
\]
Дискриминант $D=1+64=65$, поэтому
\[
\boxed{x=\frac{-1+\sqrt{65}}2}\quad\text{или}\quad
\boxed{x=\frac{-1-\sqrt{65}}2}.
\]

\textit{Случай II: $t=\dfrac35$.}
\[
\frac{x+3}{5}-\frac{2}{x+3}=\frac35.
\]
Умножая на $5(x+3)$, получаем
\[
(x+3)^2-10=3(x+3),
\]
\[
x^2+6x+9-10-3x-9=0,
\]
\[
x^2+3x-10=0,
\]
\[
(x+5)(x-2)=0.
\]
Значит,
\[
\boxed{x=-5}\quad\text{или}\quad\boxed{x=2}.
\]

\textbf{Шаг 5. Отбираем корни на отрезке $[-6;-4]$.}
\begin{itemize}[leftmargin=*,itemsep=0.25em]
\item Число $-5$ принадлежит отрезку.
\item Число $2$ не принадлежит отрезку.
\item Корень $\dfrac{-1+\sqrt{65}}2$ положителен, поэтому не подходит.
\item Поскольку $8<\sqrt{65}<9$, получаем
\[
-5<\frac{-1-\sqrt{65}}{2}<-4{,}5.
\]
Этот корень принадлежит отрезку $[-6;-4]$.
\end{itemize}

\subblock{Проверка}
Все четыре найденных числа отличны от $-3$, поэтому удовлетворяют ОДЗ. Для каждого корня выражение
\[
t=\frac{x+3}{5}-\frac{2}{x+3}
\]
принимает одно из значений $1$ и $\dfrac35$. Подстановка в преобразованное исходное равенство даёт соответственно
\[
5\cdot1^2+4=8\cdot1+1=9
\]
и
\[
5\cdot\left(\frac35\right)^2+4
=\frac{29}{5}
=8\cdot\frac35+1.
\]
Так как все переходы при $x\ne-3$ равносильны, найденные числа являются корнями исходного уравнения.

\subblock{Итог}
\result для пункта \textbf{а)} получены все корни:
\[
\boxed{x=-5,\quad x=2,\quad
x=\frac{-1-\sqrt{65}}2,\quad
x=\frac{-1+\sqrt{65}}2.}
\]
Для пункта \textbf{б)} на отрезке $[-6;-4]$ лежат только
\[
\boxed{x=-5,\qquad x=\frac{-1-\sqrt{65}}2.}
\]
Основной навык для закрепления --- узнавать повторяющееся выражение и сохранять равносильность уравнения при каждом преобразовании.

\subblock{Задача на закрепление}
\textbf{а)} Решите уравнение
\[
\frac{(x-1)^2}{4}+\frac{16}{(x-1)^2}
=6\left(\frac{x-1}{4}-\frac{2}{x-1}\right)+2.
\]
\textbf{б)} Укажите его корни, принадлежащие отрезку $[-5;0]$.

\end{document}